\documentclass[12pt]{article}
\usepackage{geometry}
\geometry{letterpaper, margin=1in}
\usepackage{hyperref}
\usepackage[usenames,dvipsnames]{xcolor}
\usepackage{microtype}

\definecolor{accent}{RGB}{26, 60, 110}   % matches resume navy

\begin{document}

%----------HEADER----------
\begin{center}
  {\LARGE\textbf{\textcolor{accent}{SITONG LI}}} \\[4pt]
  \small Minneapolis, MN
  \;\textbullet\;
  \href{mailto:ruc.sitongli@gmail.com}{ruc.sitongli@gmail.com}
  \;\textbullet\;
  \href{https://linkedin.com/in/YOURHANDLE}{LinkedIn}   % <-- fill in
  \;\textbullet\;
  \href{https://github.com/YOURHANDLE}{GitHub}          % <-- fill in
  \;\textbullet\;
  (952) 201-4026
\end{center}

\vspace{1em}

\noindent \today

\vspace{1em}

\noindent Dear Professor Dube,

\vspace{1em}

I am writing to apply for the Research Professional position in AI and Machine
Learning for Social Impact at the Becker Friedman Institute for Economics. My
long-term goal is to pursue a Ph.D.\ at the intersection of machine learning and
causal inference, and your research program---applying NLP to body-worn camera
footage to study and improve policing---represents exactly the kind of work I want
to build toward.

Your 2025 \textit{Quarterly Journal of Economics} paper, ``A Cognitive View of
Policing,'' struck me as a genuinely novel contribution: rather than treating
officer behavior as a black box, it attempts to recover latent mental states from
observable actions. The methodological challenge is hard in both directions---the
NLP pipeline must be rigorous enough to yield reliable signal from noisy, high-volume
audio-visual data, and the causal identification must be clean enough to support
inference about what actually changes officer behavior. That combination is rarer
than it should be in applied work, and it is the combination I most want to develop.
The prospect of working on the field experiment arm---evaluating whether
AI-powered tools actually change how officers interact with the public---adds
a further dimension I find compelling: building the measurement infrastructure is
only half the problem; the other half is designing an evaluation that can
credibly answer the question.

I bring strong technical foundations for this work. I am proficient in Python and
have built production-quality pipelines for processing large unstructured corpora,
including text and audio data. I have hands-on experience with the Hugging Face
ecosystem---fine-tuning transformer models for classification, sequence labeling,
and information extraction---and with scikit-learn for structured-data workflows.
I have also worked extensively with large language models for document-level
information extraction, with careful attention to prompt robustness and output
validation. My experience is not limited to the modeling side: I have built the
full pipeline, from raw data ingestion and cleaning through feature engineering,
model training, evaluation, and reproducible reporting. On the causal side, I have
hands-on experience with difference-in-differences and event-study designs, and I
am actively developing proficiency in Stata for the field experiment work.

Before coming to the United States, I worked as a fixed-income analyst at Harvest
Fund Management covering local government bonds in China. That experience gave me a
firsthand sense of how institutional behavior---by governments, regulators, and
enforcement agencies---shapes outcomes for ordinary people, and why measuring that
behavior rigorously matters. It also taught me to work with messy, high-dimensional
administrative data under real time pressure. Those habits have carried directly into
my research work.

I am strongly committed to the full two-year term and am available to begin
July 1, 2026. I would welcome the chance to discuss how my technical background
fits your current projects.

\vspace{1em}

\noindent Sincerely,\\[0.5em]
\textbf{Sitong Li}

\end{document}
